%% kapa: exam=1 round=28 year=2017 period=1 form=A subject=economics subject_ko=경제학원론 seq=5 number=45
%% source: Q-Net 감정평가사 기출문제 (cst00309) artlSeq=5199603 file=extracted/1008/5199603/2017년 제28회 감정평가사 1차시험 문제지/1교시 A형.pdf
%% conversion: kordoc 4.24.51 -> tools/gen1.py (auto); review=proofread-sonnet; flags=none
\begin{question}{45}
동일한 콥-더글러스(Cobb-Douglas) 효용함수를 갖는 甲과 乙이 X재와 Y재를 소비한다. 다음 조건에 부합하는 설명으로 옳지 않은 것은?
\begin{framedtext}
○ 초기에 甲은 X재 10단위와 Y재 10단위를 가지고 있으며, 乙은 X재 10단위와 Y재 20단위를 가지고 있다. \par
○ 두 사람이 파레토(Pareto) 효율성이 달성되는 자원배분상태에 도달하는 교환을 한다.
\end{framedtext}
\begin{choices}
  \item 교환 후 甲은 X재보다 Y재를 더 많이 소비하게 된다.
  \item 교환 후 甲은 X재와 Y재를 3 : 5의 비율로 소비하게 된다.
  \item 교환 후 乙은 X재를 10 단위 이상 소비하게 된다.
  \item 교환 후 두 소비자가 각각 Y재를 15 단위씩 소비하는 경우는 발생하지 않는다.
  \item 계약곡선(contract curve)은 직선의 형태를 갖는다.
\end{choices}
\end{question}
\answer{2}
